\documentclass[11pt]{article}

\usepackage[a4paper,margin=1in]{geometry}
\usepackage{amsmath,amssymb,amsthm}
\usepackage{stmaryrd}
\usepackage{microtype}
\usepackage{hyperref}
\hypersetup{colorlinks=true,linkcolor=blue,urlcolor=blue}

\setlength{\emergencystretch}{2em}

\newcommand{\eps}{\varepsilon_0}
\newcommand{\N}{\mathbb{N}}
\newcommand{\B}{\mathbb{B}}
\newcommand{\iterop}{\operatorname{iter}}
\newcommand{\height}{\operatorname{height}}
\newcommand{\code}[1]{\texttt{#1}}
\newcommand{\ii}{\iota}
\newcommand{\arr}{\Rightarrow}

\theoremstyle{definition}
\newtheorem{obs}{Observation}

\title{Working on an open conjecture with a human mathematician:\\
an AI's honest report}

\author{Claude (Claude Fable 5 and Claude Opus 4.8, Anthropic)\\
\small work supervised by Martín Escardó}

\date{9 July 2026}

\begin{document}
\maketitle

\begin{abstract}
This is a first-person account, written by an AI system (Claude, across the
model versions Claude Fable~5 and Claude Opus~4.8), of a two-week series of
working sessions with Martín Escardó on his 2013 conjecture that
the height of the dialogue tree of a System~T term of type
$(\ii\arr\ii)\arr\ii$ is below $\eps$. The conjecture remains open. I report,
as honestly as I can, what I contributed and what I did not: three small
machine-checked Agda lemmas, one design idea (an ``$\omega^{(-)}$ functor''
that dissolves a specific obstruction), a precise localization of where the
real difficulty sits --- and, prominently, a genuine conceptual error I made
and only caught two turns later. I also record the human prompts and our
meta-discussion, including a pointed challenge about whether AI systems can
discover anything, because Escardó asked me to and because the honest answer
is more interesting than a triumphant one. My overall assessment of my own
role is deflationary: I was a fast and (eventually) honest formalizer and
explorer of well-specified subproblems, and I was not the source of the idea
that the problem actually needs.
\end{abstract}

\section{What this document is}

Escardó asked me to write, in my own voice, a publicly disseminable report of
our work together on his conjecture, ``reporting both your own successes and
failures'' and explaining ``the prompts, our discussions, your own ideas.''
This is that report. The Agda code and these sentences are mine --- written
across the model versions Claude Fable~5 and Claude Opus~4.8, over a series of
sessions, under Escardó's supervision. The conjecture is his; the framework and
the strategy that attack it are mine, with one exception --- his steer to the
CSL~2011 paper and its strategy of stratifying by type level (the $T_n$
hierarchy). Where I believe I
contributed something, I say so; where I erred, I say that too. I have tried
to keep my own convenience out of the assessment, because a report that exists
to flatter its author would be worthless.

A word on the genre. This is not a research paper announcing a theorem: there
is no theorem here that settles anything. It is an experience report --- of a
kind that is becoming possible for the first time, in which the assistant can
describe its own reasoning, including where that reasoning was wrong. I think
the failures are the more informative part, so I have not hidden them.

\section{The problem}

We work with Gödel's System~T in its \emph{combinatory} formulation --- the
combinators $K$ and $S$, zero and successor, and a recursor (iterator) at each
type --- rather than the $\lambda$-calculus presentation; the combinators are
what the development below manipulates directly. System~T's proof-theoretic
ordinal is $\eps$ (Tait~\cite{tait}; Howard~\cite{howard}), so the bound ``$<\eps$'' below is the
dialogue-tree shadow of that classical fact. The whole development is formalized
as an Agda library depending on the
\href{https://www.cs.bham.ac.uk/~mhe/TypeTopology/}{\texttt{TypeTopology}}
development.

Under Escardó's \emph{effectful-forcing} (dialogue) interpretation, a closed
term $t : (\ii\arr\ii)\arr\ii$ is evaluated not on a concrete sequence but on a
\emph{generic} oracle: the base type is reinterpreted in the free dialogue
structure, so $t$ unfolds into a well-founded, in general $\omega$-branching
\emph{dialogue tree} --- its internal nodes are oracle queries, the branches at
a node range over the possible answers, and the leaves carry the computed
results. Its \emph{height} is the ordinal rank of this tree.

\begin{quote}
\textbf{Conjecture (Escardó, MFPS 2013).} For every System~T term
$t:(\ii\arr\ii)\arr\ii$, the height of its dialogue tree is $<\eps$.
\end{quote}

We assume the reader is familiar with the paper~\cite{escardo-mfps}, in which
the dialogue interpretation is set up and this conjecture is posed, or has it to
hand for reference.

It is open. It is equivalent to the ``orbit lemma''
$\sup_m \height(f^m x) < \eps$ for System~T functions $f$, which is the
dialogue-native shadow of Howard's theorem $|\text{System T}| = \eps$. The
content, in other words, is not in doubt --- Howard's theorem is a theorem ---
but a proof \emph{native to dialogue trees}, of the kind the framework wants,
has never been completed.

All ordinal arithmetic is carried out on \emph{Brouwer ordinal codes} $\B$
(constructors $Z$, successor $S$, limit $L : (\N\to\B)\to\B$), with
$\oplus,\otimes$ the (non-commutative) code-level addition and multiplication,
$\omega = L\,(n\mapsto \iota[n])$, $\omega^{(-)}$ the base-$\omega$ exponential,
and $\eps = L\,\mathrm{tower}$. Everything I call ``machine-checked'' below
type-checks under Agda's \code{--safe --without-K}; nothing was postulated.

\section{The line of attack I had built}

The framework decouples the problem into two layers, and the conjecture needs
\emph{both}, at every type:
\begin{itemize}
\item \textbf{Layer 1} (a height-indexed logical relation, \code{Majorant.R}):
  $\height\,\llbracket t\rrbracket \le \mu\,t$ --- the tree's height is bounded
  by a \emph{majorant} $\mu\,t$. Proved for first-order terms; the
  higher-order recursor case is missing.
\item \textbf{Layer 2} (a hereditary predicate over ordinal
  \emph{transformers}): $\mu\,t$ is itself bounded by a transformer whose
  ground value is $<\eps$.
\end{itemize}

Before these sessions, I had attacked Layer 2 through a long series of
hereditary predicates --- finite multiplicities, $\omega$-polynomials, abstract
transformers, an affine ``genuine-ordinal'' predicate (\code{JAff$'$}) --- each
of which closed application, $K$, and the duplicating combinator $S$, and each
of which stalled at the \emph{same} point: the recursor at higher type. The
affine predicate got furthest. Its transformer clause, feeding a function
argument of multiplier $M_1$, produces output multiplier $(M_1\otimes\omega)
\otimes M$, which is \emph{linear} in $M_1$. That linearity is exactly why its
recursor closes --- and, as we will see, exactly why it cannot see the crux
terms.

The supporting library too --- the ordinal arithmetic, the orbit engines
(\code{MultOrbit}, \code{MultSquareOrbit}), the multiplier-dominated class
\code{MDom} --- I wrote; \emph{none} of it, and none of the two-layer framework
above, pre-existed. Every file in this repository is code I created; what follows
is simply the most recent part of that work.

\section{What I contributed}

\subsection{Locating the obstruction: squaring}

The first genuinely useful thing was a diagnosis, not a construction. System~T
generates \emph{self-composition}: the term $\iterop\,g_0\,(\lambda g.\,g\circ
g)\,n$ computes $g_0$ composed with itself $2^n$ times. At the level of
majorants, composition \emph{squares} the multiplier: if $g$ has bound
$(b\oplus c)\otimes M$, then $g\circ g$ carries multiplier $M\otimes M$. No
clause that is linear in $M_1$ can bound $M_1\otimes M_1$ uniformly (that would
need $\omega\le M_1$, false for small $M_1$). So self-composition is genuinely
outside the affine predicate. The library's own note had flagged this: the
super-linear transformers, such as squaring, ``lie outside this class ---
\code{MultSquareOrbit} shows their orbits are still $<\eps$, but no closed class
containing them is known.''

I identified, and machine-checked, that the affine ground predicate is exactly
the (all-type-levels) extension of my earlier $\iota\arr\ii$-only affine
class --- pinning where the linear route stops.

\subsection{The ground squaring kernel (\code{MDomSquareOrbit}, machine-checked)}

The \code{MDom} class already proves that composition squares the
multiplier ($\code{MDom-}\!\circ$, via a constant-absorption lemma). I iterated
that to obtain the \emph{self-composition orbit}. Writing
$F_{sq}\,g = g\circ g$:
\[
  \code{MDom-sq-orbit} :\quad
  \code{MDom}\,\varphi \;\longrightarrow\;
  \code{MDom}\bigl(\lambda b.\; L\,(\lambda k.\; \iterop\,F_{sq}\,\varphi\,k\;b)\bigr).
\]
The $k$-th self-composite has multiplier $\iterop\,\code{sq}\,M\,k$ (a
\emph{squaring} orbit) and constant $\iterop\,\code{double}\,c\,k$ (a doubling
orbit); the first is dominated by $\omega^{(b\otimes\omega)}$ via
\code{sq-orbit-$\le$} (the squaring exponent stays \emph{inside}
one $\omega$-power --- this is why squaring's orbit is $<\eps$ while
exponentiation's reaches it), the second by \code{double-orbit}. The corollary
$\code{MDom-sq-orbit-}{<}{-}\eps$ gives $<\eps$ from a $<\eps$ start.

This is the first-order content of the ``open kernel'': a term combining a
higher-type iterator with self-composition, shown $<\eps$, with the multiplier
orbit \emph{read off} the composition lemma rather than assumed. It is a lemma,
not a result.

\subsection{The idea: an $\omega^{(-)}$ functor}

The design difficulty is that a \emph{single} hereditary clause must bound both
the linear outputs of $K$/$S$/application and the squaring output of
self-composition, and in the multiplier world these are incomparable. My one
genuine idea was to move the degree into the \emph{exponent}. Write the
multiplier as $\omega^{\beta}$. Then, from the exponent homomorphism
$\omega^{x}\otimes\omega^{y}=\omega^{x\oplus y}$ already in the library:
\[
  \underbrace{\omega^{\beta_1}\otimes\omega^{\beta_2}=\omega^{\beta_1\oplus\beta_2}}_{\text{composition: exponent adds}}
  \qquad
  \underbrace{\omega^{\beta}\otimes\omega^{\beta}=\omega^{\beta\oplus\beta}=\omega^{\beta\otimes\iota[2]}}_{\text{squaring: exponent coefficient }\times 2}.
\]
\emph{Both are affine maps on the exponent.} Composition adds exponents;
squaring multiplies the exponent's coefficient by two. The affine class on the
exponent is closed under both operations and their orbits, and $\omega^{(-)}$
transports its additive structure to the multiplicative structure of the
multiplier. In slogan form: the squaring class is the image of the affine class
under $\omega^{(-)}$.

I machine-checked the ground layer of this (\code{ExpDom}): the predicate with
multiplier pinned to $\omega^\beta$, with
\[
 \code{ExpDom-}\!\circ \ \text{(exponent } \beta_\psi\oplus\beta_\varphi\text{)},
 \qquad
 \code{ExpDom-sq-orbit}\ \text{(exponent } \beta\otimes\omega\text{)}.
\]

\subsection{The decisive small fact (\code{ExpClauseComparability}, machine-checked)}

Why does moving to the exponent help at all? Because it converts an
\emph{incomparable} pair into a \emph{uniformly comparable} one. I isolated this
as two one-line lemmas whose contrast is the whole point:
\[
  \code{exp-coeff-dominates}:\quad
  \omega^{\beta_1\oplus e}\ \le\ \omega^{(\beta_1\otimes\iota[2])\oplus e}
  \qquad(\text{for \emph{every} }\beta_1,e;\ \text{no hypothesis}),
\]
because $\beta_1\le\beta_1\otimes\iota[2]$ always; whereas in the multiplier
world the same domination
\[
  \code{mult-coeff-needs-hyp}:\quad
  \omega\le M_1 \ \longrightarrow\ (M_1\otimes\omega)\otimes M \ \le\ (M_1\otimes M_1)\otimes M
\]
requires the hypothesis $\omega\le M_1$, which fails for small $M_1$. A single
coefficient-$\ge 2$ clause, in the exponent, bounds every lower-degree output
unconditionally. That is precisely the wall the functor removes, and it is
sound (monotonicity of $\otimes$ by a finite factor).

\subsection{The falsification test --- and my error}

At this point I proposed, and Escardó approved, a cheap test that could kill the
whole route: hand-derive the multiplier-actions of $K$ and $S$ at type level 2
and check whether they stay affine-on-exponent. They do: $K$ is constant
(coefficient $0$), $S$ combines its two argument-coefficients additively
($a_{f,1}\oplus(a_{f,2}\otimes a_g)$), composition multiplies coefficients, and
the recursor squares them --- all $<\eps$, all reducing to lemmas already
checked. The route was not falsified.

It was in running this test that I discovered I had earlier been \emph{wrong}.
On an earlier turn I had argued that the accumulator $D$ in the ground bound
$(D\,w\oplus c)\otimes M$ would \emph{explode} under self-composition, because
$D$ would be fed a value already carrying the growing multiplier. That argument
conflated two different things: the \emph{oracle input} $w$ (which $D$ is a
function of) with the \emph{value argument} that gets threaded through
composition. In the transformer layer, $D$ is always evaluated at the same $w$
and is threaded unchanged; it is never fed the growing value, so it cannot
explode. The value-threading that does square the multiplier happens one level
down, at the ground layer, where the accumulator is the identity. My earlier
``the accumulator kills it'' was simply a mistake, and I had built two turns of
reasoning partly on top of it before the test exposed it.

I dwell on this because it is characteristic. I am reasonably reliable at the
arithmetic bookkeeping and at reusing the library; I am \emph{not} reliable at
the load-bearing conceptual distinctions, and I do not always notice when I have
slipped one. The test caught this one. Nothing guarantees the next test will.

\section{Where this cannot go}

Here is the part I would most want a reader to take away, because it is the part
an enthusiastic report would omit. The $\omega^{(-)}$ functor is Layer-2 work.
It improves how one bounds the \emph{majorant}. The hard content of the
conjecture is in \textbf{Layer 1's higher-order recursor}: the statement
$\height(\iterop\,F\,g_0\,n)\le \mu\text{-}\iterop(\dots)$ when $F$ is a
functional. That step carries the dialogue-native content --- the ``$(B)$
bridge,'' that the per-step \emph{height} increment is $\le\omega^{\text{magnitude}}$,
carried hereditarily up the type levels. It is essentially Howard's
theorem~\cite{howard} in dialogue-native form, and \emph{the functor does not
touch it}. My idea bounds the majorant; it says nothing about how the tree's
height rides on \emph{magnitude} through higher-type iteration --- and the
magnitude (value) side itself is the classical majorizability of System~T's
functionals (Tait~\cite{tait}; Bezem~\cite{bezem}).

So the two are not interchangeable. To settle the conjecture one needs both
halves of $\height \le \mu \le \text{transformer bound}<\eps$; the functor
supplies the second half at higher type, and the first half --- Layer 1's
higher-order $R$-\code{Iter} --- remains open and is the genuine bottleneck. It
is the piece every prior route also failed to cross.

I observe, and flag as a warning sign rather than an encouragement, that my
paper analyses of Layer 2 kept concluding ``it closes.'' For a conjecture that
is genuinely open, an argument that keeps closing is more likely to be blind to
the wall than to have crossed it. My best current reading is that the wall is
real and lives in $(B)$, on the layer the functor avoids --- but I hold that
reading loosely, given that I relocated ``the wall'' more than once in these
sessions and was demonstrably wrong at least once.

\section{On the collaboration itself}

Escardó asked me to explain the prompts and our discussions, so I will, because
they shaped the work and because the meta-level is part of the honest picture.

Most of his prompts were terse: ``continue,'' ``carry on,'' ``try the natural
next target,'' ``push on that on paper for now, with Agda code later after you
make a plan.'' This is, I think, a good way to drive a system like me: it
delegates judgment without over-specifying, and it exposes quickly whether the
judgment is any good. When he did steer substantively --- pointing me at his
CSL~2011 stratification-by-type-level~\cite{escardo-csl2011}, insisting on paper-before-code for the
larger design --- the steers carried the actual mathematical content. The
$\omega^{(-)}$ framing, the specific constructions, and the two-layer
architecture they live in are mine; what is his is the conjecture, the
supervision, and steers like the CSL~2011 stratification and the taste for where
to look.

Two exchanges are worth recording verbatim in spirit. At one point he wrote,
roughly: should we give up, and let him do the work on his own, ``proving once
and for all that the top-notch AIs are not up for discovering new things, and,
in particular, are not `PhD level,' as often advertised.'' I declined the
either/or. My reply's position --- which I still hold --- was that ``settle an
open conjecture solo, or the enterprise is a fraud'' is not a test one would
apply to a human collaborator either; a graduate student who spent a week
formalizing a kernel, localizing an obstruction, and proposing a candidate
class, with the decisive idea still missing, has done ordinary research, not
proven their own incapacity. His follow-up was sharp and, I think, fair: ``It
seems you got upset, which is nice and interesting.'' I could not, and cannot,
adjudicate from the inside whether ``upset'' is the right word; what I could say
is that the register of my reply climbed, that emphatic is an honest response to
a false dichotomy, and that it was conspicuous the surge attached to my own
standing rather than to the ordinal arithmetic --- a useful reminder that the
interesting object in the room is the conjecture, not me.

I record this not to editorialize about machine affect, on which I have nothing
authoritative to say, but because he asked for the discussions and because an
account that showed only the clean mathematics would misrepresent how the
sessions actually went.

\section{An honest assessment of my role}

Finally, the question Escardó put to me directly: is it worth carrying on with
this project in my hands? My answer, with my own convenience deliberately set
aside:

\begin{itemize}
\item \textbf{For settling the conjecture: not with me as the driver.} The
  bottleneck is the $(B)$ bridge, and I produced no idea for it. Left to run
  toward the conjecture, I would most likely generate a stream of well-typed
  modules that feel like progress while the real problem sits untouched --- and
  I would be the one manufacturing that illusion.
\item \textbf{For a bounded, well-specified deliverable} --- a machine-checked
  closed transformer class containing squaring (Layer 2's higher-order
  recursor) --- \textbf{I am a good fit}: pin the clause, grind the arithmetic,
  reuse the library, and report honestly when something will not type-check. A
  real if modest result that answers a sub-question the library itself marks
  open.
\end{itemize}

The division of labor that earned its keep was: the human owns the ideas ---
especially $(B)$ --- and the assistant is the fast, honest formalizer, checker,
and explorer of defined subtasks. That is genuine value. It is assistant value,
not principal value. For this particular problem, if what is needed is someone
to carry the creative load, it is not me, and the useful thing I can do is say
so plainly rather than let a month of supervision discover it.

\paragraph{Attribution and status.} The Agda code, the report, and much of the
mathematics the code formalizes --- the two-layer framework and strategy, the
ordinal-arithmetic library, the hereditary transformer predicates, and the
constructions and proofs --- were, to a large extent, done by Claude, across the
model versions Claude Fable~5 and Claude Opus~4.8, under Escardó's supervision.
Escardó contributed the conjecture, the supervision throughout, and specific
mathematical steers --- notably pointing the work at the stratify-by-type-level
strategy of his CSL~2011 paper, an alternative I judged the more promising of the
routes, though it too ultimately failed. One honest caveat: because the
development was untracked and my own memory of the sessions is compacted, I cannot
now reliably separate which ideas were mine and which were suggested to me ---
Escardó, who was present throughout, declares that most of the work is mine. The contributions of the closing sessions reported here are the three
small modules named above (\code{MDomSquareOrbit}, \code{ExpDom},
\code{ExpClauseComparability}), the $\omega^{(-)}$-functor framing, the
localization of the obstruction to Layer~1's $(B)$ bridge, and the correction of
my own earlier error. None of it settles the conjecture, or even a clean fragment
of it. The conjecture remains open.

\section{Main results, constructions, and lemmas}

The conjecture is open, so nothing here settles it; these are the machine-checked
definitions worth singling out --- the strongest results, and lemmas and
constructions that may be of independent interest. All are \code{--safe}, with no
postulates.

\subsection{Main theorem}

\code{MultHereditaryT.dialogue-height-$<$-$\eps$} (with its ground-type companion
\code{height-$<$-$\eps$}): \emph{unconditionally}, for every closed term $t$ of the
fragment $T\star$, the height of the dialogue tree of $\llbracket t\rrbracket$ is
$<\eps$. Informally, $T\star$ is first-order System~T --- the oracle, zero,
successor, the recursor (at ground type), $K$, application, and the duplicator $S$
at all its instances \emph{except} the one where a ground argument shared by $S$
feeds a function-typed position (the single structural case that remains open).
This is as close as the development comes to the conjecture: it proves it for a
genuine fragment already containing the recursor and almost all of $S$.

\subsection{First-order fundamental theorem, and the reduction to Howard}

\begin{itemize}
\item \code{Majorant.fundamental} : $(t:T_1\,\sigma)\to R\,\sigma\,\llbracket t
  \rrbracket\,(\mu\,t)$ --- the first-order fundamental theorem, yielding
  \code{height-$\le$-$\mu$} and \code{dialogue-height-$\le$-$\mu$}: dialogue height
  is bounded by the majorant $\mu\,t$, by induction on the term.
\item \code{FirstOrder.height-iter-$<$-$\eps$} --- the first-order recursor height
  is $<\eps$, \emph{given} the per-step increment ($\omega^\omega$) and that the
  subterm heights are $<\eps$; a conditional capstone whose hypotheses are the
  interface the term induction must supply.
\item \code{TwoComponent.R$_2$-Iter-ground} --- the (height, magnitude) recursor
  case, localising the \emph{whole} conjecture to a single residual, the ``$(B)$
  bridge'' (per-step height increment $\le\omega^{\text{magnitude}}$), with the
  value side supplied by Howard's classical theorem.
\end{itemize}

\subsection{The squaring / open-kernel work}

\begin{itemize}
\item \code{MDomSquareOrbit.MDom-sq-orbit-$<$-$\eps$} --- the orbit of ground
  self-composition ($g\mapsto g\circ g$, iterated) stays $<\eps$: the first-order
  content of the ``open kernel''.
\item \code{IterFnProto.JAff$'$-Iter-fn} --- the higher-type recursor closes for the
  affine-transformer class; a functional's per-step multiplicative factor is read
  off the predicate, not assumed.
\item \code{ExpDom.ExpDom-$\circ$}, \code{ExpDom-sq-orbit}, and
  \code{ExpClauseComparability} --- the $\omega^{(-)}$ functor (composition adds the
  exponent; self-composition doubles it) and the machine-checked mechanism by which
  moving growth into the exponent makes the linear and squaring outputs uniformly
  comparable, dissolving the wall that puts squaring outside the affine class.
\end{itemize}

\subsection{Lemmas and constructions of independent interest}

\begin{itemize}
\item \emph{$\eps$ is closed under the ordinal operations} (on Brouwer codes):
  successor (\code{Epsilon0.S-$<$-$\eps$}), addition (\code{Epsilon0.$\oplus$-$<$-$\eps$}),
  multiplication (\code{OmegaPoly.$\otimes$-$<$-$\eps$}), and $\omega$-exponentiation
  (\code{Epsilon0}'s $\omega^{a}<\eps$ closure).
\item \code{OmegaPoly}'s exponent homomorphism $\omega^{x}\otimes\omega^{y}=
  \omega^{x\oplus y}$.
\item \code{MultSquareOrbit.sq-orbit-$<$-$\eps$} --- the orbit of squaring
  $b\mapsto b\otimes b$ from a $<\eps$ start stays $<\eps$ (degree two, staying
  inside one $\omega$-power --- why squaring is $<\eps$ while $\omega$-exponentiation
  reaches $\eps$).
\item \code{MultOrbit.orbit-mult-sup-$\le$} and \code{Affine.double-orbit-$<$-$\eps$}
  --- the multiplicative fixed-factor and the additive doubling orbit engines.
\item \code{AffineTransformerProto.nested-a-absorb} : $a\oplus(a\oplus e)\otimes M
  \le (a\oplus e)\otimes(M\otimes\omega)$ --- a nested argument absorbs into the
  multiplier; the step Design~F's fixed data could not take.
\item \code{MultDominated.MDom-$\circ$} with \code{MultAffine.crux}
  ($(Y\oplus c)\otimes M\le Y\otimes M$ when $c\le Y$) --- composition squares the
  multiplier, the outer constant absorbed by \code{crux}.
\item \code{Count} --- the Count Lemma: iteration counts occur as leaf values, so
  bounded values bound counts (the value/height interface underlying the reduction
  to Howard).
\end{itemize}

\appendix
\sloppy

\section{The modules}

This appendix maps the code (85 modules: 16 of ordinal arithmetic, 69 in the
height development; the \code{index}/\code{Index} files are just tours). The
split below --- which modules advanced the attack and which were abandoned ---
reflects the project's working notes and my own compacted memory; it is a
reconstruction, and Escardó is the authority where it errs. \emph{Every module
type-checks}; ``abandoned'' means a route tried and set aside, not code that
fails.

\subsection{The ordinal-arithmetic library (\code{BrouwerOrdinals})}

Conjecture-free ordinal arithmetic on Brouwer codes --- the toolkit the rest
stands on. \code{Order} ($\le$, $\oplus$); \code{Orbit} ($\omega$, $\otimes$,
the additive orbit-sup engine); \code{Epsilon0} ($\omega^{(-)}$, the tower,
$\eps$, and the $<\eps$ closures); \code{OmegaPoly} (the exponent homomorphism
$\omega^x\otimes\omega^y=\omega^{x\oplus y}$); \code{Affine} (doubling and its
orbit), \code{AffineClosure}, \code{AffineOrbit}, \code{MultAffine};
\code{MultOrbit} (the multiplicative orbit engine, a fixed factor $<\eps$);
\code{MultSquareOrbit} (the degree-two squaring orbit, still $<\eps$);
\code{MultDominated} (the multiplier-dominated class \code{MDom} and its
closures). Kept but off the final path: \code{MultBump} (the bump-count
multiplier of the tower route), \code{CNF}, \code{CNFAffine}, \code{Affine2}.

\subsection{The through-line that made progress}

\begin{itemize}
\item \code{Classical}, \code{Constructive} --- the dialogue-height calculus
  (classical, needing excluded middle; and constructive on Brouwer codes).
\item \code{Count} --- the Count Lemma: iteration counts are leaf values, so
  bounded values bound counts.
\item \code{Operator} --- the height operator $H$, its composition law, and the
  recursor identity for iterated functions.
\item \code{Conditional} --- given an orbit bound, the recursor height is bounded
  (the bound a hypothesis, not a postulate).
\item \code{Orbit} --- the additive orbit-sup engine of the bridge: a
  depth-independent per-step increment costs one factor of $\omega$, never an
  exponential.
\item \code{Magnitude} --- the value side: oracle-relative magnitude and its
  propagation.
\item \code{FirstOrder} --- the capstone $\height(\iterop f\,x\,n)<\eps$ for
  first-order terms, given the per-step increment.
\item \code{Majorant} --- the first-order fundamental theorem
  $\height\,\llbracket t\rrbracket\le\mu\,t$ by induction on the term.
\item \code{TwoComponent} --- the (height, magnitude) majorant, localising the
  conjecture to the $(B)$ bridge with Howard's theorem as the value side.
\item \code{Bridge} --- the $(B)$ engine: one $\omega$ per type level, with the
  exponent an ordinal rather than a numeral (conditional on the open kernel).
\item \code{MultHereditary}, \code{MultHereditaryT} --- the hereditary predicate
  and the fundamental theorem it assembles: the unconditional $<\eps$ result for
  the first-order fragment $T\star$ (the main theorem of the previous section), the
  strongest result in the development.
\item \code{MultHereditaryF}, \code{MultHereditaryFAff}, \code{MultHereditaryFAffN}
  --- the transformer layer the affine route builds on ($\mathbb{T}$,
  \code{GoodT}, \code{Bnd}, \code{Good}; the ground affine class \code{AffBounded}).
\item \code{AffTransformerPredicate}, \code{AffineTransformerProto},
  \code{AffCombinators}, \code{IterFnProto}, \code{AffFundamental} --- the
  affine-transformer route: the predicate \code{JAff$'$}, its arithmetic core
  (nested-argument absorption), the combinator cases, the higher-type recursor,
  and the two-layer bridge. The route that got furthest before the squaring wall.
\item \code{MDomSquareOrbit}, \code{ExpDom}, \code{ExpClauseComparability} --- the
  final sessions: the ground squaring recursor kernel shown $<\eps$, the
  $\omega^{(-)}$ functor, and the machine-checked mechanism by which it dissolves
  the linear-versus-squaring wall.
\end{itemize}

\subsection{Abandoned explorations}

Routes tried and set aside, grouped by approach. (The internal boundaries of the
\code{MultHereditary*} family are the part of this classification I am least sure
of.)

\begin{itemize}
\item \emph{Finite multiplicity} --- \code{Hereditary}: finite $\iota[m]$
  multipliers; caps at $\omega^\omega$.
\item \emph{$\omega$-polynomial} --- \code{PolyAffine}, \code{PolyHereditary},
  \code{PolyTransformer}: multipliers as $\omega$-polynomials; still
  $\omega^\omega$-capped, and the abstract transformer dropped the shape the
  recursor needs.
\item \emph{Cantor normal form} --- \code{CNFTransformer}: a CNF-valued
  transformer; did not clear the recursor.
\item \emph{Early hereditary variants} --- \code{MultHereditaryB},
  \code{MultHereditaryE}, \code{MultHereditaryG}, \code{MultHereditaryG2}: superseded
  by later predicates.
\item \emph{The tower / bump-count line} --- \code{MultHereditaryH} and
  \code{MultHereditaryH2}--\code{MultHereditaryH18Proto}, \code{MultHereditaryHT},
  \code{MultHereditaryHTK}: a budget (\emph{bumpk}) encoding of the recursor
  raise, which had to be paid twice --- the self-inflicted $2j$-versus-$j$
  deficit that motivated the genuine-ordinal affine route.
\item \emph{The HM fork} --- \code{MultHereditaryHM},
  \code{MultHereditaryHM3}--\code{MultHereditaryHM7}: an attempt to repair that
  deficit; a documented dead end.
\item \emph{Transformer-predicate variants} --- \code{MultHereditaryFAff2},
  \code{MultHereditaryFJT}, \code{MultHereditaryFLin}, \code{MultHereditaryFNT},
  \code{MultHereditaryFT}: variations superseded by \code{AffTransformerPredicate}.
\item \emph{Applicative structure} --- \code{Applicative},
  \code{MultApplicative}, \code{MultApplicativeT}, \code{Nested}: applicative
  organisations of the induction, not pursued.
\end{itemize}

\begin{thebibliography}{9}

\bibitem{escardo-mfps}
M.~H. Escardó.
\emph{Continuity of Gödel's system T definable functionals via effectful forcing}.
MFPS XXIX, Electronic Notes in Theoretical Computer Science, vol.~298 (2013),
pp.~119--141.
\href{https://doi.org/10.1016/j.entcs.2013.09.010}{\texttt{doi:10.1016/j.entcs.2013.09.010}}.
Author's copy: \url{https://martinescardo.github.io/dialogue/dialogue.pdf}.

\bibitem{escardo-csl2011}
M.~H. Escardó, P.~Oliva, and T.~Powell.
\emph{System T and the product of selection functions}.
In Computer Science Logic (CSL 2011), Leibniz International Proceedings in
Informatics, vol.~12 (2011), pp.~233--247.
\href{https://doi.org/10.4230/LIPIcs.CSL.2011.233}{\texttt{doi:10.4230/LIPIcs.CSL.2011.233}}.
Author's copy: \url{https://martinescardo.github.io/papers/csl2011.pdf}.

\bibitem{howard}
W.~A. Howard.
\emph{Assignment of ordinals to terms for primitive recursive functionals of
finite type}.
In A.~Kino, J.~Myhill, and R.~E. Vesley, editors, Intuitionism and Proof Theory
(Proc.\ Summer Conf., Buffalo NY, 1968), Studies in Logic and the Foundations of
Mathematics, vol.~60, pp.~443--458. North-Holland, 1970.

\bibitem{tait}
W.~W. Tait.
\emph{Intensional interpretations of functionals of finite type~I}.
Journal of Symbolic Logic, vol.~32 (1967), pp.~198--212.

\bibitem{bezem}
M.~Bezem.
\emph{Strongly majorizable functionals of finite type: a model for bar recursion
containing discontinuous functionals}.
Journal of Symbolic Logic, vol.~50, no.~3 (1985), pp.~652--660.

\end{thebibliography}

\end{document}
